\section{Experimental Setup}
\label{sec:setup}

\subsection{Dataset and Settings}
\label{sec:setup-dataset}
We use the dataset introduced in Heo et al.~\cite{heo2025study} for our experiments. The dataset has a total of 6,600 issue reports gathered from eleven active open-source projects:
\texttt{ansible/ansible}, \texttt{bitcoin/bitcoin},
\texttt{dart-lang/sdk}, \texttt{dotnet/roslyn},
\texttt{facebook/react}, \texttt{flutter/flutter},
\texttt{kubernetes/kubernetes}, \texttt{microsoft/TypeScript},
\texttt{microsoft/vscode}, \texttt{opencv/opencv}, and
\texttt{tensorflow/tensorflow}. Each project has 600 issue reports
with an equal distribution across the labels
(\textit{bug, feature, question}). To support their LLM fine-tuning
method, Heo et al. divided the dataset into 50/50 train and test
splits, stratified by project and label. This results in 300 train
and 300 test issue reports per project.

We also adopt Heo et al.'s project-specific (PS) and project-agnostic
(PA) configurations to examine how data scope affects IRC performance.
In the PS setting, each project's train and test data are used in
isolation, resulting in eleven different train-test pairs. In the PA
setting, the train splits of all projects are merged into a single
training set of 3,300 issue reports. In both configurations, \ragtag
and \votag only use the training splits as their retrieval index to
ensure fair comparison with the fine-tuning baseline.

\subsection{Models}
\label{sec:setup-models}
We use the Qwen2.5-Instruct model family~\cite{yang2024qwen25} in our evaluations with four different parameter sizes: 3B, 7B, 14B, and 32B. Qwen2.5 is an open-source LLM with strong performance on general-purpose benchmarks. Using the same model at different parameter sizes isolates the effect of model scale from model architecture across LLM-based IRC approaches. Model temperature is set at 0.1 for deterministic classifications.


\subsection{Evaluation Metrics}
\label{sec:setup-metrics}

We report macro-averaged precision, recall, and $F_1$ over the three labels (\texttt{bug}, \texttt{feature}, \texttt{question}), along with overall accuracy. Per-class metrics are reported when the asymmetry across classes is itself the subject of analysis. Outputs that fail to parse a valid label (after the XML-tag extraction and regex fallback described in \Cref{sec:approach}) are counted as incorrect, and an \emph{invalid rate} is reported separately so any parsing-induced loss is auditable.

\paragraph{Aggregation convention (PA vs PS).}
The PA setting produces a single set of predictions over the full 3{,}300-issue test set, and we compute macro $F_1$ directly over that set. The PS setting produces eleven sets of predictions, one per project, each over its own ${\sim}300$-issue test slice. Throughout the paper we report PS results using \emph{pooled} aggregation: the eleven per-project prediction sets are concatenated into one 3{,}300-issue test set and macro $F_1$ is computed once on the pool. This treats each test issue as one observation in both settings and yields apples-to-apples comparisons between PA and PS. The alternative --- computing macro $F_1$ separately within each project and averaging the eleven numbers --- weights projects rather than issues; because macro $F_1$ is non-linear in the underlying confusion-matrix counts, the two conventions disagree, and on our data the per-project mean is consistently 0.002--0.011 macro $F_1$ below the pooled value with the gap largest for fine-tuning. We adopt pooling as the canonical reporting convention to remove this asymmetry; per-project breakdowns appear only where the project is the unit of analysis and are explicitly labelled as such.

\subsection{Setup Hardware}
\label{sec:setup-hardware}
% Points to cover:
% 24. All experiments were run on a server with the following specifications: NVIDIA RTX A6000 GPU (48 GB VRAM), [CPU model + core count], [RAM amount], [OS / CUDA version].
% 25. TODO: confirm exact hardware specs (see paper/TODO.md) -- need to check via kubectl on the cluster.
